\documentclass[10pt,a4paper]{amsart}
\usepackage[margin=19mm]{geometry}
\usepackage{amsmath,amssymb,mathtools}
\usepackage[scr=rsfs,cal=euler]{mathalfa}
\usepackage{xfrac}
\usepackage[unicode,hidelinks,bookmarks=false]{hyperref}
\newcommand{\ie}{i.e.\ }
\newcommand{\cf}{cf.\ }
\newcommand{\resp}{respectively\ }
\newcommand{\Subsection}{\subsection}
\providecommand{\nobreakdash}{\nobreak\hbox{-}}
\setlength{\emergencystretch}{2em}
\multlinegap0pt
\usepackage{amssymb}
\usepackage{tikz}
\usepackage{enumerate}
\newtheorem{proposition}{Proposition}
\title{Quasi-pseudometric modular spaces as $\mathscr{Q}$-categories}
\author{C\'esar L\'opez-Pastor, Tatiana Pedraza, Jes\'us Rodr\'{\i}guez-L\'opez}
\date{Source: arXiv:2604.27588v1 (CC BY 4.0)}
\begin{document}
\maketitle

\noindent\emph{Source and licence.} Quasi-pseudometric modular spaces as $\mathscr{Q}$-categories. Version arXiv:2604.27588v1 (CC BY 4.0), distributed by arXiv under the Creative Commons Attribution 4.0 International licence, http://creativecommons.org/licenses/by/4.0/, as recorded in the arXiv licence metadata for this exact version and shown on the abstract page https://arxiv.org/abs/2604.27588v1. Full licence notice: \texttt{licenses/arxiv-2604.27588-CC-BY-4.0.txt}.\par\smallskip

\noindent\textit{Editorial scope.} Counted unit: the article's proposition prop:civalueq (source0.tex lines 865--870). The article's complete original proof of it is delivered with it (source0.tex lines 872--910, one \texttt{proof} environment of 39 lines). No mathematics is written, repaired or re-derived: the statement and the proof are the article's own lines, in the article's own notation, and no retained line is substituted or re-flowed.  No further source-local context is needed: every cross-reference of every retained line resolves inside the two delivered spans.  Nothing was omitted from any retained span, so the package carries no omission record.  No retained line contains a citation, so the wrapper carries no bibliography and no bibliography entry.  No retained line contains a figure environment, an image-inclusion command or drawing code, so no image is delivered and the package declares no asset dependency.  Editorial formatting only, in addition: an AMS \texttt{amsart} preamble that restates the article's own declarations and macros used by the retained lines, namely the environments \texttt{proposition} on separate counters, together with the standard packages those lines need.  The retained environments print in this document as Proposition 1 in order of appearance, whereas the article prints the same items with the numbering of its own full document.  The complete native source of the article as distributed in the arXiv e-print for this exact version is bundled unchanged as \texttt{sources/199485/source0.tex}.
	\begin{proposition}\label{prop:civalueq}
		Consider the binary operation $\oplus:\nabla\times\nabla\to\nabla$ given by
		\[(f\oplus g)(t):=\bigvee_{r+s\leq t}^{\leq^\text{\normalfont op}}(f(r)+g(s))=\bigvee_{r+s=t}^{\leq^\text{\normalfont op}}(f(r)+g(s))\]
		for all $t>0.$ 
		Then $(\nabla,\leq^\text{\normalfont op},\oplus)$ and $(\nabla_L,\leq^\text{\normalfont op},\oplus)$ are  CI-value quantales.
	\end{proposition}

	\begin{proof}
		
		It is straightforward to check that $f\oplus g\in\nabla$ for every $f,g\in\nabla.$
		
		
		Given $t>0,$ we prove that
		$$\bigvee_{r+s=t}^{\leq^\text{\normalfont op}}f(r)+g(s)=\bigvee_{r+s\leq t}^{\leq^\text{\normalfont op}}f(r)+g(s).$$
		
		
		Notice that $\{f(r)+ g(s): r+s=t\}\subseteq \{f(r)+g(s): r+s\leq t\}$, so
		\[\bigvee_{r+s=t}^{\leq^\text{\normalfont op}}(f(r)+g(s))\leq\bigvee_{r+s\leq t}^{\leq^\text{\normalfont op}}(f(r)+ g(s))=(f\oplus g)(t).\]
		On the other hand, if $r+s\leq t$, then $s\leq t-r$. By isotonicity of $g$, $g(s)\leq^\text{\normalfont op} g(t-r)$ so $$f(r)+g(s)\leq^\text{\normalfont op} f(r)+g(t-r)\leq^\text{\normalfont op}\bigvee_{r\leq t}^{\leq^\text{\normalfont op}}(f(r)+ g(t-r))=\bigvee_{r+s=t}^{\leq^\text{\normalfont op}}(f(r)+g(s)).$$ Hence,
		\[(f\oplus g)(t)=\bigvee_{r+s\leq t}^{\leq^\text{\normalfont op}}(f(r)+g(s))\leq^\text{\normalfont op} \bigvee_{r+s=t}^{\leq^\text{\normalfont op}}(f(r)+g(s)),\]
		proving the desired equality.
		
		We next check that $(\nabla,\oplus)$ is a commutative monoid.
		
		First, commutativity is clear from the commutativity of the sum. Furthermore, given $f\in\nabla$,
		\[(f\oplus\mathbf{0})(t)=\bigvee_{r+s=t}^{\leq^\text{\normalfont op}}f(r)+\mathbf{0}(s)=\bigvee_{r+s=t}^{\leq^\text{\normalfont op}}f(r)=f(t),\]
		where the last inequality holds since $f$ is isotone. So $\mathbf{0}$ is the neutral element for $\oplus.$
		
		
		Finally, to prove the associativity property, one has that
		\begin{align*}
			((f\oplus g)\oplus h)(t)&=\bigvee_{r+s=t}^{\leq^\text{\normalfont op}}(f\oplus g)(r)+h(s)=\bigvee_{r+s=t}^{\leq^\text{\normalfont op}}\left(\bigvee_{u+v=r}f(u)+g(v)\right)+h(s)=\\			&=\bigvee_{\substack{r+s=t \\ u+v=r}}^{\leq^\text{\normalfont op}}f(u)+g(v)+ h(s)=\bigvee_{u+v+s=t}^{\leq^\text{\normalfont op}}f(u)+ g(v)+ h(s).
		\end{align*}
		where the last equalities hold by the continuity of the sum. Since the final expression does not depend on the order of the elements, we can assure that $\oplus$ is associative.
		
		
		
		To prove the distributivity property of $\oplus$ with respect to suprema, we only need to show it for one side because of the commutativity of the operation. For any $\{g_i\}_{i\in I}\subseteq\nabla$,
		\begin{align*}
			\left(f\oplus \bigvee_{i\in I}^{\leq^\text{\normalfont op}}g_i\right)(t)&=\bigvee_{r+s=t}^{\leq^\text{\normalfont op}}\left(f(r)+\bigvee_{i\in I}g_i(s)\right)=\bigvee_{r+s=t}^{\leq^\text{\normalfont op}}\bigvee_{i\in I}^{\leq^\text{\normalfont op}}f(r)+ g_i(s)\\&=\bigvee_{i\in I}^{\leq^\text{\normalfont op}}\bigvee_{r+s=t}^{\leq^\text{\normalfont op}}f(r)+ g_i(s)=\left(\bigvee_{i\in I}(f\oplus g_i)\right)(t).
		\end{align*}
		
		Finally, since $f\leq^\text{\normalfont op}\mathbf{0}$ for all $f\in\nabla$ then $\mathbf{0}=\top$ so the quantale $(\nabla,\leq^\text{\normalfont op},\oplus)$ is integral.
		
		Additionally, a routine check shows that $f\oplus g\in\nabla_L$ for every $f,g\in\nabla_L.$ Since $(\nabla_L,\leq^\text{\normalfont op})$ is a sublattice of $(\nabla,\leq^\text{\normalfont op})$ with the same top and bottom, then $(\nabla_L,\leq^\text{\normalfont op},\oplus)$ is also a CI-quantale.
	\end{proof}

\end{document}
